\subsection{Pancasila Sebagai Ideologi Negara}
\noindent \textbf{Idea: }Sesuatu yang ada dipikiran sebagai hasil rumusan ide, gagasan, konsep, pengertian dasar.


\noindent \textbf{Ideologi: }Ilmu pengetahuan tentang ide-ide/gagasan tetang pengertian-pengertian dasar.


\noindent \textbf{Ideologi muncul pertama kali: }Masa revolusi Prancis diperkenalkan oleh Destutt De Tracy (abad 18 dan 19).


\noindent \textbf{Logos: }Ilmu pengetahuan atau teori.


Jika pancasila negara dihubungkan dengan tujuan negara maka: \textit{``Pancasila sebagai ideologi negara''}.


\subsubsection{Contoh Ideologi Sebagai Nilai Cita-Cita}
Pancasila tidak bersifat konkret melainkan luwes(fleksibel).


\noindent{Kedudukan pancasila sebagai pandangan hidup:}
\begin{itemize}
\item Kehidupan berbangsa: Relasi nonformal $\rightarrow$ norma.
\item Kehidupan bernegara: Formal $\rightarrow$ Dasar negara $\rightarrow$ ideologi.
\end{itemize}


\subsubsection{Ciri-Ciri Ideologi Terbuka:}
\begin{enumerate}
\item Nilai yang berasal dari cita-cita masyarakat.
\item Hasil musyawarah.
\item Bersifat dinamis.
\end{enumerate}

\noindent \textbf{Ciri-Ciri Pancasila Sebagai Ideologi Terbuka:}
\begin{enumerate}
\item Nilai dan cita-cita tidak dipaksakan dari luar.
\item Ideologi melalui musyawarah/konsensus.
\end{enumerate}

\subsubsection{Fungsi Ideologi}
\noindent \textbf{Menurut Yudi Latif (2021):}
\begin{enumerate}
\item Menawarkan kerangka tertib sosial.
\item Mengembangkan model masa depan yang dihendaki.
\item Menjelaskan transformasi politik.
\end{enumerate}

\noindent \textbf{Peran Pancasila Sebagai Ideologi:}
\begin{itemize}
\item Pedoman.
\item Penyaring nilai-nilai yang tidak sesuai.
\end{itemize}


\noindent \textbf{Pengamalan Pancasila yang Tidak Sesuai:}


Ketetapan MPR XVII/MPR/1998 $\rightarrow$ Tujuan dan cita-cita ideologi.
\begin{enumerate}
\item Struktur kognitif dan orientasi dasar $\rightarrow$ orientasi dasar.
\item Norma menjadi pedoman.
\item Sumber motivasi/kekuatan mencapai tujuan.
\item Sarana pendidikan.
\end{enumerate}
